\section*{Chapter \ref{ch:ll:where-latency-comes-from} --- Where Latency Comes From}

\begin{solution}{exo:ll:where-latency-comes-from:1}
$45 \times 1.462 / (2.998\times10^8) = \qty{219}{\nano\second}$ each way.
\end{solution}

\begin{solution}{exo:ll:where-latency-comes-from:2}
$148 \times 8 = 1\,184$ bits: \qty{118.4}{\nano\second} at 10~Gb/s and \qty{47.4}{\nano\second} at 25~Gb/s.
\end{solution}

\begin{solution}{exo:ll:where-latency-comes-from:3}
Means add: \qty{1\,030}{\nano\second}. Nothing exact can be said about the median of the sum from the stage means or
medians; for light-tailed, independent stages it is usually close to the sum of the medians, as in the chapter's
synthetic path (\qty{2.17}{\micro\second} against \qty{2.13}{\micro\second}).
\end{solution}

\begin{solution}{exo:ll:where-latency-comes-from:4}
$1 - 0.996^5 = 1.98\%$. Since more than 1\% of messages meet a stall, the path's p99 lies in the stalls, although no
stage's p99 does.
\end{solution}

\begin{solution}{exo:ll:where-latency-comes-from:5}
By \cref{prop:ll:where-latency-comes-from:union} with $\alpha_i = 0.001/4$: each stage at its 99.975th percentile, with
targets summing to at most \qty{10}{\micro\second}.
\end{solution}

\begin{solution}{exo:ll:where-latency-comes-from:6}
At \qty{2.0}{\micro\second}, removing the stalls lifts $\P(\text{win})$ from 0.662 to 0.675 (1.3 points). Along the
curve with stalls, $\P(\text{win})$ rises by about 0.81 per microsecond near \qty{2}{\micro\second}, so the same gain
takes a median about \qty{16}{\nano\second} lower.
\end{solution}

\begin{solution}{exo:ll:where-latency-comes-from:7}
\texttt{ll\_budget.composition(stall\_p=0.05)}: each stage's p99 now lies in its stalls, and the sum of the six stage
p99s is \qty{79}{\micro\second}, against an end-to-end p99 of \qty{32}{\micro\second}. The comparison has flipped: the
sum of percentiles understates when stalls are rarer than the percentile, and overstates when they are common and do
not coincide. Neither is a budget.
\end{solution}

\begin{solution}{exo:ll:where-latency-comes-from:8}
Percentiles do not add. The three p99s were measured in separate runs, so they say nothing about how the stages'
slow messages line up: rare stalls that each benchmark excludes from its p99 can put the pipeline's p99 far above
\qty{1.3}{\micro\second} (\cref{ex:ll:where-latency-comes-from:nonsub}), and stalls common to all three can put it
below. Measure the pipeline end to end, message by message, with timestamps at each stage boundary.
\end{solution}

\begin{solution}{pb:ll:where-latency-comes-from:1}
\begin{enumerate}
\item $\P(X_A < X_B)$, with ties shared; it assumes our \omterm{def:ll:where-latency-comes-from:latency}{latency} and the competitor's on a given message are
independent (a burst that slows both breaks this).
\item $\ln X_A - \ln X_B \sim \mathcal N(\ln(2.0/2.2), 2\times0.15^2)$, so $\P(\text{win}) =
\Phi\big(\ln(2.2/2.0)/(0.15\sqrt2)\big) = \Phi(0.449) = 0.673$.
\item 66.1\%.
\item We lose only the stalled races we would otherwise have won, about $0.02 \times 0.67$; and the competitor's own
stalls hand us a few.
\item p50 \qty{2.01}{\micro\second}; p99 \qty{23.6}{\micro\second}.
\item Halving the median: 98.0\%. Halving the stall rate: 66.7\%, with the p99 falling to \qty{3.5}{\micro\second}.
\item 67.4\%: 1.3 points.
\item The median.
\item Stalls start at $10^5 \times 0.02 = 2\,000$ a second and last \qty{30}{\micro\second}: 6\% of the time.
\item $20 \times 0.06 = 1.2$ stale fills a second at $0.5 \times \$0.01 \times 100 = \$0.50$: \$0.60 a second,
\$14\,040 a day.
\item It halves it, to \$7\,020 a day.
\item Exposure depends on the tail (stall rate times stall length); the race on the body.
\item About \qty{1.64}{\micro\second}.
\item $1.64 - 1.0 = \qty{0.64}{\micro\second}$ over five stages: about \qty{128}{\nano\second} each.
\item At the 99.8th percentile ($1 - 0.01/5$).
\item \textbf{Named result.} We win 66.1\% of races; halving our median lifts it to 98.0\%, halving our stall rate only
to 66.7\%, although it cuts our p99 from \qty{23.6}{\micro\second} to \qty{3.5}{\micro\second}.
\item Races depend on the whole body of both distributions and exposure on the tail; one number is neither.
\item Its \omterm{def:ll:where-latency-comes-from:latency}{latency} is visible only through outcomes: the timestamps of its orders against ours in the exchange's
message data, or the share of races lost as a function of our own measured \omterm{def:ll:where-latency-comes-from:latency}{latency}.
\item A schedule-driven execution algorithm working a large order over hours.
\item A shared, falsifiable definition of ``fast enough'' for every stage, so that each change is measured against it.
\end{enumerate}
\end{solution}

\begin{solution}{iq:ll:where-latency-comes-from:1}
Tick-to-trade runs from the triggering market-data message to the resulting order; wire-to-wire is the same interval
taken on the cable, by a capture device with hardware timestamps on a tap or switch, so it includes the network cards
and the operating system. Inside the program, timestamp the packet's arrival and the order's hand-off with the
processor's clock.
\iqlookfor{the two instants named explicitly, and who takes each timestamp.}
\end{solution}

\begin{solution}{iq:ll:where-latency-comes-from:2}
Losses concentrate in the slow messages (stale quotes, missed races) and the mean hides them; high percentiles show
the tail, and the median the ordinary message.
\iqlookfor{that \omterm{def:ll:where-latency-comes-from:latency}{latency} is a distribution and that money is lost in its tail.}
\end{solution}

\begin{solution}{iq:ll:where-latency-comes-from:3}
Not much: percentiles do not add. The union bound gives $q_{0.97} \le \qty{6}{\micro\second}$ for any dependence, but
the pipeline's p99 can exceed \qty{6}{\micro\second} if each component has rare stalls below its own 1\% level. Measure
it end to end.
\iqlookfor{the non-additivity of quantiles and the union bound.}
\end{solution}

\begin{solution}{iq:ll:where-latency-comes-from:4}
The network card's receive and transmit paths, interrupts and the kernel's network stack (or a kernel-bypass
library), copying, decoding, the book update, the decision, risk checks, encoding, and waiting: for the thread to be
scheduled, for a lock, for memory that is not in cache.
\iqlookfor{hardware, operating system and software stages, and the stalls between them.}
\end{solution}

\begin{solution}{iq:ll:where-latency-comes-from:5}
Races are decided by the body of both distributions: at 60\% the median is where the return is, unless measurement
shows that lost races coincide with our stalls. The tail deserves money for another reason: exposure of resting quotes
while stale.
\iqlookfor{the distinction between race outcomes and exposure, and a request for data before spending.}
\end{solution}

\begin{solution}{iq:ll:where-latency-comes-from:6}
Capture packets in and out on a tap with hardware timestamps for the wire-to-wire figure; add timestamps with the
processor's clock at every stage boundary inside the program, logged off the \omterm{def:ll:where-latency-comes-from:hot}{hot path}; join the two by message
identifiers; replay a recorded busy day at its original pace; report per stage and end to end at p50, p99 and p99.9
from joint per-message samples.
\iqlookfor{external and internal timestamps, joined per message, under a realistic load.}
\end{solution}
